\documentclass[11pt]{article}
\usepackage[margin=1in]{geometry}
\usepackage{amsmath,amssymb,amsthm,mathtools}
\usepackage[hidelinks]{hyperref}
\newtheorem{theorem}{Theorem}
\newtheorem{lemma}[theorem]{Lemma}
\newtheorem{corollary}[theorem]{Corollary}
\newcommand{\Q}{\mathbb Q}
\newcommand{\F}{\mathbb F}
\title{Uniform local Newton lemmas for factorial moment polynomials}
\author{User-supplied quantitative argument checked and integrated; no priority claim}
\date{16 November 2026}

\begin{document}
\maketitle
For positive integers $n,m$, write
\[
 F_{n,m}(t)=\sum_{k=0}^n\binom nk(mk)!t^k.
\]
The following local statements are the inputs used by the
\href{ALL_UNIVARIATE_TWO_TERM_SFC.tex}{all-index constant-plus-power theorem}.
They do not leave an exceptional set in that completed ordinary family.

\section{Exact local statements for every exponent}
\begin{lemma}\label{lem:shift}
Let $n=ap+i$ with $a\ge1$, $i\ge0$ and prime $p>m\max(a,i)$.
For $j=bp+c$, $0\le c<p$, $0\le j\le n$,
\[
 v_p\!\left(\binom nj(mj)!\right)
 =mb+\left\lfloor\frac{mc}{p}\right\rfloor
       +\mathbf1_{\{c>i\}}. \tag{1}
\]
The Newton polygon in increasing coefficient degree has a horizontal edge
of length $i$ (absent if $i=0$), followed by an edge of width $ap$ and
slope $m/p$, of reduced denominator $p$. Moreover
\[
 \overline{F_{n,m}}=\overline{F_{i,m}}\quad\text{in }\F_p[T]. \tag{2}
\]
For a rational square divisor $g^2\mid F_{n,m}$,
\[
 \deg g=u+vp,\qquad u,v\ge0,\quad2u\le i,\quad2v\le a. \tag{3}
\]
If $\deg g<p$, a primitive integral normalization has degree-preserving
reduction on the horizontal edge.
\end{lemma}
\begin{proof}
We have $a,i<p$ and
$mn=map+mi\le(p-1)p+(p-1)<p^2$. Thus Legendre has only its first term,
$\lfloor mj/p\rfloor=mb+\lfloor mc/p\rfloor$.
Kummer contributes one borrow precisely when $c>i$, with no second borrow
because $j\le n<p^2$. This proves (1).
For $j\le i$ the valuation is zero. Subtract $m(j-i)/p$ from (1).
If $c\le i$, then $mc\le mi<p$ and the difference is
$m(i-c)/p\ge0$, with equality exactly when $c=i$.
If $c>i$, the difference is
$1-\{mc/p\}+mi/p>0$.
The endpoints $j=i,n$ therefore establish the stated sloped edge.
Since $p>m$, its reduced denominator is $p$.

For $j\le i$, Lucas gives $\binom nj\equiv\binom ij\pmod p$.
For $j>i$, (1) is positive: either $b\ge1$, or $b=0,c>i$.
This proves the full identity (2).
Newton multiplication allocates doubled horizontal width $2u$ and doubled
sloped width $2vp$ to $g^2$, proving (3).
By Gauss's lemma choose $g$ primitive integral with integer quotient in
$F_{n,m}$. The constant coefficient one forces $g(0)=\pm1$.
When its degree is below $p$, it has only horizontal slopes, hence leading
$p$-adic valuation zero and degree-preserving reduction.
No real-root parity condition is used.
\end{proof}

\begin{corollary}\label{cor:preceding}
Write $n=q_n+d$, where $q_n$ is the largest prime at most $n$.
If $q_n>m\max(1,d)$, every rational square divisor has degree at most
$\lfloor d/2\rfloor$. In particular this applies when
$d<\sqrt n$ and $n>(m+1)^2$.
\end{corollary}
\begin{proof}
Apply Lemma~\ref{lem:shift} with $a=1$, $p=q_n$ and $i=d$.
Then $2v\le1$ forces $v=0$. For the final sufficient condition,
$q_n=n-d>n-\sqrt n>m\sqrt n\ge m\max(1,d)$.
\end{proof}


\end{document}
